\chapter*{Conclusion Générale}
\phantomsection
\addcontentsline{toc}{chapter}{Conclusion Générale}

Le projet \textbf{eduGenius} a permis de concevoir et de réaliser une plateforme éducative innovante intégrant les dernières avancées en intelligence artificielle générative. Au cours de ce travail, nous avons pu relever les défis liés à la gestion du temps réel, au stockage cloud et à l'interaction avec des modèles de langage complexes dans un environnement agile.

Les objectifs fixés au début du projet ont été atteints, notamment la création d'un écosystème fluide pour les enseignants et les étudiants. La méthodologie Scrum a prouvé son efficacité pour piloter un projet aux technologies changeantes, permettant de transformer chaque besoin en un incrément fonctionnel et testé.

Cependant, ce projet reste ouvert à des améliorations futures. Nous envisageons d'intégrer un tutorat personnalisé par l'IA capable de répondre aux questions des étudiants en temps réel pendant les sessions live, ainsi que d'approfondir l'analyse prédictive pour identifier précocement les étudiants en difficulté.

Ce stage a été une expérience enrichissante, me permettant de consolider mes compétences en développement Full-stack (MERN) et d'explorer le potentiel de l'IA dans le secteur de l'éducation.
